%%%%%%%%%%%%%%%%%%%%%%%%%%%%%%%%%%%%%%%%%%%%%%%%%%%%%%%%%%%%%%%%%%%%%%%%%
%                                                                       %
%                        Chapter 8 - Section 1                          % 
%                                                                       %
%%%%%%%%%%%%%%%%%%%%%%%%%%%%%%%%%%%%%%%%%%%%%%%%%%%%%%%%%%%%%%%%%%%%%%%%%

\chapter{Dynamical Systems}

A recurring theme in this book is the use of mathematical models
consisting of a set of differential equations to explore the behavior
of physical systems as they evolve over time.  Some examples we have
encountered are the $S$-$I$-$R$ epidemiological model, predator-prey
systems, and the motion of a pendulum. We call such a set of
differential equations a {\bf dynamical system}\index{dynamical
  system}. Dynamical systems play important roles in all branches of
science.  In this chapter we will develop some general tools for
thinking about them, with particular emphasis on the kinds of
geometric insight provided by the concepts of {\bf state space} and
{\bf vector field}.

\section{State Spaces and Vector Fields}

If you look back at the examples we've considered, 
%
\mar{The standard way\\ to graph solutions}
%
many of them take the following form: we have two (or more) variable
quantities $x$ and $y$ that are functions of time, and we want to find
the nature of these functions.  What we have to work with is a model
for the way the functions $x(t)$ and $y(t)$ are changing---i.e., we
are told how to calculate $x'(t)$ and $y'(t)$ whenever we know the
values of $x$ and $y$, and possibly $t$.  From a given starting point,
we typically used something like Euler's method to get values for $x$
and $y$ at times on either side of the starting value.  We then
graphed the solutions as functions of time---$x$ against $t$ and $y$
against $t$.

In many instances, the rules determining $x'(t)$ and $y'(t)$ depend
only on the current values of $x$ and $y$, but not on the value of
$t$, so that knowing the current state of the system (as specified by
its $x$ and $y$ values) is sufficient to determine the future and past
states of the system.  Such systems are said to be {\bf
  autonomous}\index{dynamical system!autonomous}.  These are the only
systems we will be considering in this chapter.

In  autonomous systems 
%
\mar{A new way\\ to graph solutions:\\ as trajectories\\ in state space} 
%
there is another way of visualizing the solutions that can be very
powerful. Instead of plotting values of $x$ and $y$ as functions of
time, we view these values as coordinates of a point in the $x$-$y$
plane.  As the system changes, the point $(x, y)$ will trace out a
curve in this plane.  The point $(x, y)$ is called a {\bf
  state}\index{state}, and the portion of the plane corresponding to
physically possible states is called the {\bf state space}\index{state
  space} of the system. The solution curves that get traced out in
state space are called {\bf trajectories}\index{trajectory}.  By
looking at three examples, we will see how this method of analysis can
help us understand the overall behavior of a system.

There are a number of effective software packages available which can
perform efficiently all the operations we will be considering, and one
of them would probably be the most useful tool for exploring the ideas
in this chapter.  On the other hand, the basic numerical operations
are quite simple, and it is easy to modify the programs developed
earlier in the text to perform these operations as well.  For those of
you who enjoy programming, we will from time to time point out some of
these modifications. It can be instructive to implement them in your
own programs, and we urge you to do so.

\subsection{Predator--Prey Models}\index{predator--prey model}

In chapter~4.1, we looked at several models for the dynamics of a
simple system consisting of foxes ($F$) and rabbits ($R$).  Our first
model was
\begin{align*}
R' &= .1 R\left( 1- \dfrac{R}{10000}\right) - .005\, R F \quad 
                  \mbox{rabbits per month}, \\
F' &= .00004 \, RF-.04 \, F \quad \mbox{foxes per month}.
\end{align*}

When we started with the initial values $R(0)=2000$ rabbits and $F(0)=
10$ foxes, Euler's method produced the following solutions for the
first 250 months:

\newpage

\mbox{}
\vspace{180pt}

\special{psfile=../ps/calc2/pred1.ps}

Let's see how the same model looks  
%
\mar{Two ways of representing the same solution graphically}
%
when we express it in the language of state spaces.  The state space
consists of points in the $R$-$F$ plane.  For physical reasons our
state space consists only of points $(R, F)$ satisfying $R\geq 0$ and
$F\geq 0$.  That is, our state space is the first quadrant of the
$R$-$F$ plane together with the bounding portions of the $R$-axis and
the $F$-axis.We can easily modify the program used to obtain the
curves in the previous picture to plot the corresponding trajectory in
the $R$-$F$ plane. We only need change the specification of the
dimensions of the viewing window and change the {\tt plot} command to
plot points with coordinates $(R, F)$ instead of $(t, R)$ and $(t,
F)$; all the rest of the calculations are unchanged.  Here's what the
same solution looks like when we do this:

\vspace*{2.5in}
\label{`trajectory'}

\special{psfile=../ps/calc2/pred2.ps}

\newpage

You should notice several things here:
\begin{itemize}

\item The trajectory looks like a spiral, moving in towards, but never
  reaching, some point at its center.  We will see later (see
  page~\pageref{eqpt}) how to determine the coordinates of this limit
  state.

\item If we had started at any other initial state with $R>0$ and
  $F>0$, we would have gotten another spiral converging to the same
  limit (try it and see).

\item From the trajectory alone, there is no way of determining the
  time at which the system passes through the different states. In
  part, this simply emphasizes that the succession of states the
  system moves through does not depend on the initial value of $t$,
  nor does it depend on the units in which $t$ is measured---if $t$
  were measured in days or years, rather than in months, the
  trajectory would be unchanged.

\end{itemize}


\mbox{} \hfill
\noindent
\begin{minipage}{2.5in}
\quad\ If we wanted to include some information about time, one way
would be to label some points on the trajectory with the associated
time value.  If we label the points every 6 months, say, we would get
the picture at the right.  Note that the points are not uniformly
spaced along the trajectory: the spacing is \linebreak
\end{minipage}
 
\special{psfile=../ps/calc2/predtime.ps}

\vspace{-9pt}
\noindent
largest between points relatively far from the origin, where the
values of $R$ and $F$ are largest.  Moreover, the closer we come to
the limit state, the tighter the spacing becomes.

Could we have foreseen some of this behavior 
%
\mar{The differential equations indicate how the state changes}
%
by looking at the original differential equations?  Since the
differential equations give $R'$ and $F'$ as functions of $R$ and $F$
alone, for each point $(R, F)$ in the state space we can calculate the
associated values for $R'$ and $F'$.  Knowing these values, we can in
turn tell in what direction and with what speed a trajectory would be
moving as it passed through the point $(R, F)$.  Using our (by now)
standard argument, in time $\Delta t$ the change in $R$ would be
$\approx R'\Delta t$, while $F$ would change by $\approx F'\Delta t$.
We can convey this information graphically by choosing a number of
points in the state space, and from each point $(R, F)$ drawing an
arrow to the point $(R+R'\Delta t, F+F'\Delta t)$. We would typically
choose a value for $\Delta t$ that keeps the arrows a reasonable size.
Here's what we get in our current example when we choose a $16\times
16$ grid of points in the region $0\leq R\leq 3000$ and $0\leq F\leq
30$, with $\Delta t =1$.

\vspace{190pt}
\label{`vecfield'}

\special{psfile=../ps/calc2/vecf1.ps}



Several things are immediately clear from this picture:  
%
\mar{Arrows indicate the change in state\ldots}
%
the arrows suggest a general counter-clockwise flow in the plane;
change is most rapid in the upper right corner; near the limit point
of the flow and near the origin change is so slow that arrows don't
even show up there.

Moreover, since the method used to construct the arrows is exactly the
way Euler's method calculates the trajectories themselves, the
solution trajectory through a given initial state is a curve
%
\mar{\ldots and trajectories are tangent to the arrows}
%
in the state space which at every point is tangent to the arrow at
that point. For instance, if we superpose the trajectory graphed on
page~\pageref{`trajectory'} on the picture above, we get the
following:

\vspace{100pt}

\special{psfile=../ps/calc2/vecf2.ps}

\newpage

\enlargethispage*{30pt}
The net effect of this construction is thus 
%
\mar{How to construct a geometrical visualization of a dynamical system}
%
to transform a problem in {\em analysis}---findinging a solution to a
system of differential equations---into a problem in {\em
  geometry}---finding a curve which is tangent everywhere to a
prescribed set of arrows. This correspondence between the analytical
and the geometrical ways of formulating a problem is very
powerful. Let's sum up the way this correspondence was established:
\begin{itemize}

\item We set up a {\bf state space} for the system being studied.
  Each point---called a {\bf state}---in the space corresponds to a
  possible pair of values the system could have.

\item There is a rule which assigns to each point in the state space a
  {\bf velocity vector}\index{velocity vector}---which can be
  visualized as an arrow in the space based at the given
  point---specifying the rates at which the coordinates of the point
  are changing.  The rule itself, which is just our original set of
  rate equations, is called a {\bf vector field}\index{vector field}.
  Geometrically, we can visualize the vector field as the state space
  with all the associated arrows.

\item Solutions to the dynamical system correspond to {\bf
  trajectories} in the state space.  At every point on a trajectory
  the associated velocity vector specified by the vector field will be
  tangent to the trajectory.  The existence and uniqueness principle
  for the solutions of differential equations---there is a unique
  solution for each set of initial values---is geometrically expressed
  by the property that every point in the state space lies on exactly
  one trajectory.  The set of all possible trajectories is called the
  {\bf phase portrait}\index{phase portrait} of the system.  For
  instance, part of the phase portrait of the system we have been
  considering appears below.  We have drawn only a few
  trajectories---if we had drawn them all, we would have seen only a
  black rectangle since there is a trajectory through every point.
\end{itemize}

% \newpage

%\mbox{} \vspace{190pt}

\vspace{130pt}

\special{psfile=../ps/calc2/phasepor.ps}

\newpage


There is almost too much detail 
%
\mar{Simplifying the picture}   
%
in the picture of the vector field and the phase portrait.  One way to
see the underlying simplicity is to notice that the space is divided
into four regions according to whether $F'$ and $R'$ are positive or
negative.  The {\em signs\/} of $F'$ and $R'$ in turn determine the
{\em direction\/} of the associated velocity vector.  For instance, if
$F'$ and $R'$ are both positive, then $F$ and $R$ must both be
increasing, which means the velocity vector will be pointing up and to
the right, while if $F'>0$ and $R'<0$, the velocity vector will be
pointing up ($F'>0$) and to the left ($R'<0$).  Let's see which states
correspond to which behaviors.  Here are original rate equations:
\begin{align*}
R' &= .1 R\left(1 - \dfrac{R}{10000}\right) - .005\,R F \quad 
          \mbox{rabbits per month}, \\
F' &= .00004\,RF-.04\,F \quad \mbox{foxes per month}.
\end{align*}

The equation for $F'$ is slightly simpler, so we'll start there.  We
see that $F'=0$ in exactly two cases:
\begin{enumerate}

\item when $F=0$, or

\item when $.00004R-.04=0$, which is equivalent to saying $R=1000$.

\end{enumerate}

The first case simply says that if we are ever on the $R$-axis
$(F=0)$, then we stay there---a trajectory starting on the $R$-axis
must move horizontally.  (If you start with no foxes, you will never
have any at a later time.)  The second case says that the value of $F$
isn't changing whenever $R=1000$.  The set of points satisfying
$R=1000$ is just a vertical line in the state space. The condition
that $F'=0$ on this line can be expressed geometrically by saying that
any trajectory crossing this line must do so horizontally (why?).

The remainder of the quadrant consists of two regions: 
%
\mar{Divide the state space into regions}
%
one consists of all points $(R, F)$ with $0\leq R<1000$ and $F>0$, the
other consists of all points $(R, F)$ with $R>1000$ and $F>0$.
Moreover, since we've already accounted for all the points where
$F'=0$, it must be true that at every point of these two regions $F' $
must be $>0$ or $<0$; $F'$ can't equal 0 in either region.  Further,
within any one region $F'$ must be always positive or always negative.
If it were positive at some points and negative at others in a single
region, there would have to be transition points where it took on the
value 0, which we have just observed can't happen.  (Be sure you see
why this is so!)  Thus to determine the sign of $F'$ in an entire
region, we only need to see what the sign is at one point in that
region.  For instance if we let $R=2000$ and $F=1$, we see that $F'=
.08-.04$, which is positive.  Therefore we will have $F'>0$ (fox
population increasing) for any other state $(R, F)$ with $R>1000$.
Similarly we can show that $F'<0$ (fox population decreasing) if
$0\leq R<1000$---the test point $R=0$ and $F=1$ is easy to evaluate.
We could, of course, have arrived at the same conclusions through more
formal algebraic arguments, which are fairly straightforward in this
instance.  In other problems, though, the ``test point'' approach may
be the more convenient.

In exactly the same way, if we look at the first rate equation, we
find that $R'=0$ in two cases:
\begin{enumerate}

\item  when $R=0$, or

\item when $.1\,(1-R/10000)-.005\,F=0$.  This is just the equation of
  a line, which can be rewritten as $F=20-.002\,R$.

\end{enumerate}

The interpretations of these two cases are similar to the preceding
analysis: any trajectory starting on the $F$-axis must stay on the
$F$-axis; any trajectory crossing the line $F=20-.002\,R$ must cross
it vertically, since $R'=0$---the $R$-value isn't
changing---there. Further, for any other state $(R, F)$ we have $R'>0$
if the point is below this line (the point $R=1$ and $F=0$ is a
convenient test point where it's easy to see without doing any
arithmetic that $R'>0$), and $R'<0$ if the point is above the line.

We can combine all this information into the following picture. We
have drawn a number of velocity vectors along the lines where $R'=0$
and $F'=0$, with one or two others in each region.

\vspace*{200pt}
\label{`equilibrium'}

\special{psfile=../ps/calc2/pred3.ps}

\newpage

We see that the entire state space is divided into four regions:
\begin{enumerate}

\item Region I, above the line $F=20-.002\,R$ and to the right of the
  line $R=1000$\,. Here $R'<0$\,, and $F'>0$\,, so all velocity
  vectors are pointing up and to the left.

\item Region II, below the line $F=20-.002\,R$ and to the right of the
  line $R=1000$\,.  Here $R'>0$\,, and $F'>0$\,, and all velocity
  vectors are pointing up and to the right.

\item Region III, below the line $F=20-.002\,R$ and to the left of the
  line $R=1000$\,.  Here $R'>0$\,, and $F'<0$\,, and all velocity
  vectors are pointing down and to the right.

\item Region IV, above the line $F=20-.002\,R$ and to the left of the
  line $R=1000$\,.  Here $R'<0$\,, and $F'<0$\,, and all velocity
  vectors are pointing down and to the left.

\end{enumerate}


Notice that this diagram makes it clear 
%
\mar{There are simple trajectories: three are just points\ldots}
%
what the limit state of the spirals is: it is the point $Q=(1000,
18)$\label{eqpt} where the line $R=1000$ and the line $F=20-.002\,R$
intersect.  Notice that at $Q$ both $R'=0$ and $F'=0$, so that if we
are ever at $Q$, we never leave---the point $Q$ is a trajectory all by
itself.  The points $O = (0, 0)$ and $P = (10\,000, 0)$ are the two
other such point trajectories. While the typical trajectory looks like
a spiral coming into the point $Q$, note that this picture contains
three other 
%
\mar{\ldots and three are straight line segments} 
%
``special'' trajectories in addition to the point trajectories:
\begin{itemize}

\item The $F$-axis for $F>0$. The point $(0, 0)$ is {\em not\/} part
  of this trajectory.

\item The portion of the $R$-axis with $0 < R < 10000$.  Here the flow
  is toward the right, towards the point $P$.

\item The portion of the $R$-axis with $10000 < R$.  Flow is to the
  left, towards $P$, with movement being slower and slower as $P$ is
  approached.  Note that this is entirely separate from the preceding
  trajectory---you can't start at any point on one of them and get to
  any point on the other.

\end{itemize}

\subsubsection{Equilibrium Points}\index{equilibrium point}

\enlargethispage*{20pt}
The three points $O$, $P$, and $Q$ 
%
\mar{Different kinds of\\ equilibrium points}
%
in the previous figure---single points which are also
trajectories---are called {\bf equilibrium points} for the system.  If
the ststem ever in such a state, it stays in it forever.  Moreover,
the system can't reach such a state from any other state (although it
may be able to come very close).  Nevertheless, the behavior of the
system is not the same near the three points. If we zoom in on each of
these points and draw some of the nearby trajectories, we get the
following pictures:

\vspace*{1.5in}

\special{psfile=../ps/calc2/eqpoints.ps}

Points $O$ and $P$ look fairly similar---they would look even more
alike if we crossed over into the negative $R$ and negative $F$
regions and included the trajectories there as well (impossible to do
in the real world, but elementary in mathematics!).  In both cases
there is one direction from which trajectories come straight towards
the point (in the case of $O$, this is the $F$-axis; for $P$ this is
the $R$-axis), and one direction in which trajectories move directly
away from the point (the $R$-axis in the case of point $O$, and the
line of slope $-.0092$ (we'll see how to find this later!) in the case
of $P$).  The remaining trajectories look sort of like hyperbolas
asymptotic to these two lines.  Equilibrium points of this sort
%
\mar{Saddle point equilibrium}
%
are called {\bf saddle points}\index{saddle point}\index{equilibrium
  point!saddle point}.  They are characterized by the property that
there is exactly one direction along which the system can be displaced
and still move back towards the equilibrium point. Displacements in
any other direction get amplified, with the state eventually moving
even further away.

Point $Q$ is quite different.  If the state experiences a small
displacement away from $Q$ in any direction, over time it will move
back towards $Q$.  Such equilibrium points are called {\bf
  attractors}\index{attractor}\index{equilibrium point!attractor},
and $Q$ is an example of a particular kind of attractor called a
%
\mar{Spiral equilibrium}
%
{\bf spiral attractor}.\index{equilibrium point!spiral
  attractor}\index{attractor!spiral} In this example, $Q$ is an
attractor for almost the entire space---if we start with any point
$(R, F)$ with $R>0$ and $F>0$, the trajectory through $(R, F)$ will
eventually come arbitrarily close to $Q$ and stay there.  We will
shortly see examples (see page~\pageref{basin}, for instance) of
attractors that draw from more limited portions of the state space.

For future reference, 
%
\mar{Attractors and repellors}  
%
we define here the concept of {\bf repellor}\index{equilibrium
  point!repellor}\index{repellor} and {\bf spiral
  repellor}\index{equilibrium point!spiral
  repellor}\index{repellor!spiral}.  Their vector fields look just
like those for the attractors, but with all the arrows reversed.  If
the state experiences a small displacement from a repellor, over time
this displacement will increase.  We will see examples of a repellor
in \ref{sec-repel} It turns out that there is a relatively small
number of kinds of equilibrium points that a system can have, and we
will meet most of them in the next several examples.  We will turn
more systematically to the problem of identifying the kinds of
equilibrium points in \ref{sec-eqs}.



\subsection{The Pendulum Revisited}\index{pendulum}\label{pend-begin}

\special{psfile=../ps/calc2/pend1a.ps}

In chapter 7 we analyzed the motion of a pendulum.  Let's see how this
analysis looks when translated into the language of state space.  We
first need to figure out what the appropriate coordinates are, which
means deciding what information we need in order to specify the state
of a pendulum.  If you look back at the model in the last chapter, you
will recall that the two variables we needed were the displacement $x$
and the velocity $v$.  Since $x$ and $v$ can potentially take on any
values, our state space will be the entire $x$-$v$ plane.  As before,
the dynamical system is specified by the equations
        $$
        x' = v, \qquad v' = -\sin{x}.
        $$
Here is what the vector field for this system looks like:

\vspace*{2.5in}

\special{psfile=../ps/calc2/pendvf.ps}

We have included in this diagram the lines where $v'=0$ (the vertical lines at every multiple of $\pi$) and the line where $x'=0$ (the horizontal line at $v=0$).  Note that the velocity vectors are horizontal on the lines corresponding to $v'=0$ and are vertical on the line corresponding to $x'=0$.
        \mar{The equilibrium points}
The points where these two sets of lines intersect---all points of the form $(k\pi, 0)$ for $k$ an integer---are the equilibrium points of the system. Let's sketch the phase portrait of this system to see more clearly what's going on:

\vspace*{2.5in}
\special{psfile=../ps/calc2/pendpha.ps}

\enlargethispage*{30pt}
We see that there are several different kinds of trajectories:
\begin{itemize}

\item There are the wavy trajectories moving from left to right across
  the top of the state space.  Note that for these trajectories the
  value of $v$ is \mar{Trajectories from left to right in the state
    plane correspond to $v>0$} always positive, and $x$ just keeps
  increasing.  These trajectories correspond to the cases where the
  velocity is great enough that the pendulum can go over the top,
  continuing to loop around counterclockwise (since $x$ is increasing
  and $x$ is measured in a counterclockwise direction) forever.
  Notice that $v$ takes on its minimum value when $x$ is an odd
  multiple of $\pi$, which is what we would expect, since the pendulum
  is at the top of its arc then.  Similarly, $v$ takes on its maximum
  value at the bottom of its arc---$x$ an even multiple of $\pi$.

\item The wavy trajectories moving from right to left across the
  bottom are similar, except that $v$ is always negative.  This
  corresponds to the pendulum wrapping around in a clockwise
  direction.

\item There are the closed loops.  Here $x$ oscillates back and forth
  \mar{Oscillations of the pendulum correspond to closed loops in the
    state plane} between some maximum and minimum value symmetrically
  placed about an even multiple of $\pi$. These trajectories
  correspond to a pendulum swinging back and forth.  The fact that
  some are centered at $x$-values other than 0 is due to the fact that
  the same position of the pendulum can be specified by an infinite
  number of values of $x$, all differing from each other by multiples
  of $2\pi$.

\item There are the equilibrium points $(k\pi, 0)$,
%
\mar{A center: a neutral equilibrium}
%
with $k$ an even integer.  This corresponds to the pendulum hanging
straight down.  If we perturb the system to a state slightly away from
such a point, the pendulum swings back and forth, and the
corresponding trajectory loops around the equilibrium point forever.
The system neither comes back to the equilibrium point---the condition
for an attractor---nor does it go wandering off even further
away---the condition for a repellor.  Such an equilibrium point is
called a {\bf center}.\index{center} Notice that it is neither an
attractor or a repellor; it is said to be a {\bf neutral
  equilibrium}.\index{equilibrium point!neutral}

\item There are the equilibrium points $(k\pi, 0)$, with $k$ an odd
  integer, correspond to the pendulum balanced vertically. These are
  saddle points---if we perturb the system slightly with {\em
    exactly\/} the right $v$-value for the given $x$-value, the system
  will move back toward the vertical position; any other combination,
  though, will cause the pendulum to wrap around and around forever or
  to oscillate back and forth forever, depending on whether the
  $v$-value is greater than or less than the critical value.

\item There are the trajectories connecting the saddle points.  These
  correspond to cases where the pendulum has just enough velocity so
  that it keeps moving closer to the vertical position without either
  overshooting and wrapping around, or coming to a stop and reversing
  direction.  In fact, these \mar{A connected curve in the phase
    portrait may be composed of more than one trajectory} trajectories
  divide the state space: on one side of such a trajectory are points
  corresponding to states where the pendulum will wrap around, and on
  the other side are points corresponding to states where the pendulum
  will swing back and forth. Note that the saddle points are {\em
    not\/} part of these trajectories, and that each arc between
  saddle points is a separate trajectory---you can't get from a point
  on one of them to a point on another.
\end{itemize}

\subsubsection{First Integrals Again}\index{first integral}

In the case of the pendulum, we have another way of thinking about the
trajectories.  Recall that in chapter~7.3 we saw that, for any
given initial conditions, the quantity $E = \tfrac{1}{2} v^2 + 1 -
\cos{x}$ was constant over time.  In the vocabulary of this chapter,
if $(x, v)$ is any state on the trajectory through $(x_1, v_1)$, then
it must be true that $\tfrac{1}{2} v^2 + 1 - \cos{x} = \tfrac{1}{2}
v_1^2 + 1 - \cos{x_1}$.  But this
%
\mar{First integrals\\ can give equations\\ of trajectories}
%
relation determines a curve in the $x$-$v$ plane.  We thus have an
algebraic condition for each of the trajectories, which will depend on
the initial values.  In this example we could actually get an equation
for each trajectory by first using the initial values to determine the
energy $E = \tfrac{1}{2} v_1^2 + 1 - \cos{x_1}$.  We could then solve
for $v$ in terms of $x$ by $v=\pm \sqrt{2(E-1+\cos{x})}$ and plot the
resulting function. (Whether we took the plus sign or the minus sign
would depend on whether the pendulum was moving counterclockwise or
clockwise.)  Let's return to our previous sketch of the phase portrait
and label some of the trajectories by their corresponding values of
$E$:

\vspace*{170pt}

\special{psfile=../ps/calc2/pendpha2.ps}

\noindent
Note that for each value of $E\geq 0$ there is more than one
trajectory having that value as its energy.

We can now characterize the different kinds of trajectories by their
associated energy $E$:
\begin{itemize}

\item If $E>2$ we get a trajectory extending from $x = \infty$ to $x=
  -\infty$ (or vice versa).

\item If $0<E<2$, the trajectory is a closed loop.

\item If $E=0$, we get a neutral equilibrium point.

\item If $E=2$, we get either a saddle point equilibrium, or a
  trajectory connecting two such saddle points.
\end{itemize}
\label{pend-end}


\subsection{A Model for the Acquisition of Immunity}\index{immunity!acquisition of}

One of the roles of mathematical 
\mar{A mathematical model can be used to think about the feasibility of a proposed explanation}
modelling is to allow researchers to explore possible mechanisms to explain an observed phenomenon.  As an example of this, consider the phenomenon of immunity:  for many infections, particularly those due to viruses, once you've been exposed to the disease your body continues to produce high levels of antibodies to the disease for the rest of your life, even in the absence of any further stimulation from the virus.   

\newpage

\aside{A capsule summary of the immune response: Vertebrates have a
  wide variety of specialized cells called {\em lymphocytes\/}
  circulating in their blood streams and lymphatic systems at all
  times.  Each lymphocyte has the ability to recognize and bind with a
  specific kind of invading organism.  The invader is called an {\em
    antigen}, and the neutralizing molecules produced by the
  responding lymphocytes are called {\em antibodies}. Prior to
  infection, the concentration level of a particular antibody is
  typically so low as to be undetectable, but the appearance of the
  antigen causes the system to respond by producing large quantities
  of the appropriate antibody.  If the body can continue to produce
  high levels of antibodies, it will be immune to reinfection.}

\enlargethispage*{30pt}
In their book {\em Infectious Diseases of Humans}, Roy
Anderson\index{Anderson, R.} \index{May, R.} and Robert May propose
the following model as a possible mechanism for how antibody levels
are sustained.  Suppose that there are two kinds of lymphocytes
(called {\em effector cells\/}) whose densities at time $t$ are
denoted by $E_1(t)$ and $E_2(t)$, with the type 2 cells being the
potential antibodies for the disease in question.  They assume further
that new cells of type $i$ ($i=1$ or $i=2$) are produced by the bone
marrow at constant rates $\Lambda_i$ and they die at a per capita
rates of $\mu_i$.  They assume that each cell type is an antigen for
the other---that is, contact with cell type 2 triggers cell type 1 to
proliferate, and vice versa.  They further assume that this
proliferation response saturates to a maximum net rate which is
dependent on the product of their respective densities.  The following
equations express this behavior:
\begin{align*}
dE_1/dt &= \Lambda_1-\mu_1E_1+a_1E_1E_2/(1+b_1E_1E_2), \\
dE_2/dt &= \Lambda_2-\mu_2E_2+a_2E_1E_2/(1+b_2E_1E_2).
\end{align*}
Here the parameters $\Lambda_i$, $\mu_i$, $a_i$, and $b_i$ would have
to be determined by experimental means.  At this stage, though, when
we are simply exploring to see if such a mechanism might account for
the phenomenon of permanent immunity, we can try a range of values for
the parameters to see how they affect the behavior of the model.

\vspace*{140pt}

\label{`amprob'}
\special{psfile=../ps/calc2/andmay1.ps}

\newpage

In the figure above, we have taken $\Lambda_1=\Lambda_2=8000$,
$\mu_1=\mu_2=1000$, $a_1=a_2=10$, and
$b_1=b_2=10^{-6}$.\label{paramvals}

\enlargethispage*{30pt}
There are a couple of features to notice about this graph:
\begin{enumerate}
        
\item Since the velocity vectors differ so much in their size,
  we \mar{Knowing the direction of the trajectories is often
    sufficient} have recorded only the direction of the velocity
  vectors, drawing all the arrows to be the same length.  Thus we
  don't really show the vector field, but its close relative, the {\bf
    direction field}\index{direction field}. This is often a useful
  substitute.

\item Since the range of values we want to represent is so
  great we have employed a common device from the sciences of plotting
  the values on a {\bf log-log scale}\index{log--log scale}.  That is,
  we have plotted the values so that each interval spanning a power of
  10---from $10^0$ to $10^1$, or from $10^3$ to $10^4$---gets the same
  space.
%
\mar{Expressing graphical information over several different orders of magnitude}
%
This is equivalent to plotting the logarithms of the values on
ordinary graph paper.  This allows us to see effects that take place
at different scales.  If we hadn't done this, but had plotted this
information on regular graph paper with the values running from 0 to
$10^5$, then some of our most interesting behavior---from $10^0$ to
$10^2$---would be compressed into the lower left-hand corner of the
graph, occupying only .001 of the vertical and horizontal scales.
\end{enumerate}

We have included in the graph the two curves corresponding to all
points satisfying $E_1'=0$ and $E_2'=0$ (note that these curves are
{\em not\/} trajectories). These curves intersect at the three points
$P_1=(8.7689, 8.7689)$, $P_2=(92.0869, 92.0869)$, and $P_3=(9907.14,
9907.14)$, which are then the equilibrium points of this system.  The
points $P_1$ and $P_3$ appear to be attractors, while the point $P_2$
is a saddle point.  In the next section (see page~\pageref{`andmay'})
we will see how to zoom in and look at the trajectories near each of
these points to confirm this impression.  Here is a picture of the
phase portrait for this system.

\vspace*{130pt}

\special{psfile=../ps/calc2/andmay3.ps}


\newpage

Note that neither $P_1$ nor $P_3$ is an attractor for the entire
system.  The {\bf basin of attraction}\index{basin of
  attraction}\label{basin} for $P_1$ appears to be a region in the
lower left of the graph, while the basin of attraction for $P_3$ is
everything else.  The boundary separating these two basins is formed
by the two heavily shaded trajectories which come toward the point
$P_2$ (since $P_2$ is a saddle point, there are only two such
trajectories---every other trajectory eventually veers off and heads
toward either $P_1$ or $P_3$).

We can now interpret this system in the following way.  State $P_1$
represents the {\bf virgin}\index{state!virgin} or {\bf resting
  state}\index{state!resting} of the system, with coordinate values
on the order of magnitude of $E_i\approx \Lambda_i /\mu_i$, which
would just be the steady state values we would have if there were no
interactions between the two kinds of cells (why?).  Note that after
small perturbations (i.e., anything roughly less than a 10-fold
increase of type 1 or type 2 cells) from $P_1$, the system will settle
back to this resting state.

Now, though, suppose a viral pathogen appears which possesses an
antigen which is identical to that expressed by cell type 1.  This has
an effect equivalent to moving vertically in the $E_1$-$E_2$ plane to
a state which is now in the basin for $P_3$. As a result, the system
immediately starts producing large quantities of type 2 cells (which
are antibodies for the virus) very rapidly, the virus is wiped out,
and the system settles into a new state---the {\bf immune
  state}\index{state!immune}---$P_3$ and remains there.  There are
now so many type 2 cells permanently floating around the body that no
further infection by the viral pathogen is possible.  The only way the
system can be switched back to state $P_1$ is if some other agent,
such as radiation therapy or infection with an HIV virus, for
instance, kills off large numbers of {\em both\/} the type 1 and type
2 cells, moving the system back into the basin of attraction for
$P_1$.  Just killing off large numbers of one type of cell won't move
the system back to state$P_1$---do you see why?

\subsection{Exercises}

\subsubsection{Two-species interactions}

We look at some variations of the predator-prey model.  While the
original context is given in terms of rabbits and foxes, similar
models can be constructed for a variety of interactions between
populations---not just predator and prey.  The key features of the
models are determined by the nature of the {\bf feedback
  structure}\index{feedback structure} between the populations.  In
the predator-prey models, the number of foxes has a negative effect on
the growth rate of rabbits---the more foxes, the slower the rabbit
population grows---while the number of rabbits has a positive effect
on the growth rate of foxes.  Can you think of other pairs of
quantities whose interaction is of this sort? In the first problem we
will look at several different models for predator-prey interactions.
In the following three problems we will look at models for other kinds
of feedback structures.

\num Below are four predator-prey models.  In each model all the
letters other than $R$ and $F$ are constant parameters.  You can
perform a general analysis, giving your answers in terms of the
unspecified parameters $a$, $b$, $c$, etc., or, if you are more
comfortable with specific values, perform the analysis using $a=.1$,
$b=.005$, $c=.00004$, $d=g=.04$, $e=.001$, $f=.05$, $h=.004$, and
$K=10,000$.  For each model you should carry out the following steps
to sketch the vector field for the model in the first quadrant of the
$R$-$F$ plane.  Compare your work with the steps that led up to the
analysis of the vector field on page~\pageref{`equilibrium'}.
\begin{itemize}

\item Write down in words a justification for each rate equation---why
  is the model a reasonable one?  What is it saying about the way
  rabbit and fox populations change?

\item Draw (in red) the set of points where $R' = 0$, and mark the
  regions where $R' > 0$ and $R' < 0$.

\item Draw (in green) the set of points where $F' = 0$, and mark the
  regions where $F' > 0$ and $F' < 0$.

\item  Mark the equilibrium points.  What color are they?

\item Sketch representative vectors of the vector field, and then
  sketch a couple of trajectories that follow these vectors.  You
  might use a computer to verify your sketches.

\item On the basis of your sketches make a conjecture about the
  stability of the equilibrium points.

\end{itemize}

\sub The original Lotka--Volterra model, proposed independently in the
mid-1920's by Lotka and Volterra\index{Lotka--Volterra model}.  This
model stimulated much of the subsequent development of mathematical
population biology.
\begin{align*}
R' &= aR-bR\,F, \\
F' &= cR\,F-dF.
\end{align*}

\sub The Leslie--Gower model.\index{Leslie--Gower model}
\begin{align*}
R' &= aR-bRF, \\
F' &= \left( e - f \dfrac{F}{R}\right) F.
\end{align*}


\sub  Leslie--Gower with carrying capacity for rabbits.
\begin{align*}
R' &= aR\left( 1- \dfrac{R}{K}\right) - b RF, \\
F' &= \left( e - f \dfrac{F}{R}\right) F.
\end{align*}

\sub  Another combination.
\begin{align*}
R' &= aR\left( 1- \dfrac{R}{K}\right) - b RF, \\
F' &= cRF+gF-hF^2.
\end{align*}

\enlargethispage*{5pt}
\num {\bf Symbiosis and mutualism}\index{symbiosis}\index{mutualism}.
Many flowers cannot pollinate themselves; instead insects like bees
transport pollen from one flower to another.  For their part, bees
collect nectar from flowers and make honey to feed new bees.  This
sort of feedback structure in which the presence of each element has a
positive effect on the growth rate of the other is called {\bf
  symbiosis} or {\bf mutualism} (there is a distinction made between
these two interactions, but mathematically they are similar).  Here is
a model: $B$ is the number of bees per acre, measured in hundreds of
bees, while $C$ is the weight of clover per acre, in thousands of
pounds.  Assume time to be measured in months.
\begin{align*}
B' & = .1(1 - .01B + .005C) B, \\
C' & = .03(1 + .04B - .1C) C.
\end{align*}

\sub Do these equations describe symbiosis?  What terms account for
symbiosis?

\sub Each equation has a negative term in it.  What aspect of reality
is this term capturing?

\sub Sketch the vector field for this system in the $B$-$C$ plane.
Find the equilibrium points, and mark them on your sketch.

\sub Draw some trajectories on your sketch, and use them to determine
the stability of the equilibrium points.

\sub Suppose an acre of land has 10,000 pounds of clover on it, and a
hive of 2,000 bees is introduced.  (What are the values of $B(0)$ and
$C(0)$ in this case?)  What happens?  Answer this question both by
drawing a trajectory and by describing the situation in words.

\sub Let a couple of years pass after the situation in part (e) has
stabilized.  Suppose the field is now mowed so only 2,000 pounds of
clover remain on it.  The bee--clover system is now at what point on
the $B$-$C$ plane?  What happens now?  Does the bee population drop?
Does it stay down, or does it recover?  Does the clover grow back?

\sub This scenario is an alternative to part (f); it is also played
out a couple of years after the situation in part (e) has stabilized.
Suppose an insecticide applied to the clover field kills two-thirds of
the bees.  The insecticide is then washed away by rain, leaving the
remaining bees unaffected.  What happens?

\num {\bf Competition}\index{competition} As a third kind of feedback
structure, consider two species X and Y competing for the same food or
territory. In this case each has a negative impact on the growth rate
of the other.  If we let $x$ and $y$ be the number of individuals of
species X and Y, respectively, then the larger $y$ is, the less
rapidly $x$ increases---and vice versa.  Here is a specific model to
consider:
\begin{align*}
x' &= .15(1 - .005x - .010y)x, \\
y' &= .03(1 -.004x - .005y)y.
\end{align*}
The term $-.010y$ in the first equation shows explicitly how an
increase in $y$ reduces the growth rate $x'$.  In the second equation
$-.004x$ tells us how much X affects the growth of Y.  Notice that Y
affects X more strongly than X affects Y.

If $x$ and $y$ are both small, then the parenthetical terms are
approximately equal to~1, so the equations reduce to
\begin{align*}
x' &= .15x, \\
y' &=  .03y.
\end{align*}
Thus, in these circumstances X's per capita growth rate is five times
as large as Y's.

In the competition for resources will the growth rate advantage permit
X to win the competition and drive out Y, or will the more adverse
effect that Y has on the growth of X permit Y to win?  Perhaps the two
species will both survive and share the resources for which they
compete.  The purpose of this exercise is to decide these questions.

\sub Suppose we start with $x = y = 10$.  What are the two growth
rates $x'$ and $y'$?  Is $x'$ about five times as large as $y'$ in
this case?  What are the approximate values of $x$ and $y$ after .5
time units have elapsed?  Is X growing significantly more rapidly than
Y?  

\sub How many equilibrium points does this system have, and where are
they?

\sub Sketch {\em and label\/} in the $x$-$y$ plane the points where
$x' = 0$ and where \mbox{$y'= 0$}.  The vector field typically points
in one of four directions: up and to the right; up and to the left;
down and to the right; or, down and to the left.  Indicate on your
sketch the zones where these different directions occur and draw
representative vectors in each zone.

\smallskip

\noindent [Note: Only three of the zones actually occur in the first
  quadrant; no vectors there point down and to the right.]

\sub Sketch on the $x$-$y$ plane the trajectory that starts at the
point $(x, y) = (10, 10)$.  Now answer the question: What happens to a
population of 10 individuals each from species X and from species Y?
In particular, does X gain an early lead?  Does X keep its lead?  Does
either X or Y eventually vanish?

\sub Is the outcome of part (d) typical, or is it not?  Try several
other starting points: $(x, y) = (150, 25)$, $(300, 10)$, $(200,
200)$, $(50, 200)$.  Do these starting points lead to the same {\em
  eventual\/} outcome, or are there different outcomes?  Use a
computer to confirm your analysis.

\sub Describe the type of each equilibrium point you found in part
(a).  Is any equilibrium an attractor?

\num \label{`basins'} {\bf Fairer competition}.  The vector field in
question 3 shows that species X didn't have a chance: all trajectories
in the first quadrant flow to the equilibrium at $(0, 200)$.  We can
attribute this to the strength of the adverse effect Y has on X---that
is, to the size of the term $-.010y$ in the first equation when
compared to the corresponding term $-.004x$ in the second equation.
Let's try to give X a better chance by increasing this term to
$-.006x$.  The equations become
\begin{align*}
x' &= .15(1 - .005x - .010y)x, \\
y' &= .03(1 -.006x - .005y)y.
\end{align*}

\sub Sketch and label in the $x$-$y$ plane the points where $x' = 0$
and where \mbox{$y' = 0$}.  Sketch representative vectors for the vector
field.  Mark all equilibrium points.

\sub What happens to a population consisting of 10 individuals each
from species X and species Y?  Is the outcome significantly different
from what it was in question 3?  To get quantitatively precise results
you will probably find a computer helpful.

\sub What happens to a population consisting of 150 individuals from
species X and 25 individuals from species Y?  Is this outcome
significantly different from what it was in question 3?

\sub Is it possible for X and Y to coexist?  What must $x$ and $y$ be?
Is that coexistence {\em stable\/}; that is, if $x$ and $y$ are
changed slightly, will the original values be restored?

\sub Sometimes X wins the competition, sometimes Y.  Mark in the
$x$-$y$ plane the dividing line between those starting points which
lead to X winning and those which lead to Y winning.

\sub Identify the type of each equilibrium point.

\sub An often-articulated concept in ecology is the {\em principle of
  competitive exclusion}\index{competitive exclusion!principle of},
which states that you can't have a stable situation in which two
species compete for the same resource---one of them will eventually
crowd out the other.  Is the model you've been exploring in this
problem consistent with such a principle?

\enlargethispage*{10pt}
\num {\bf More on the Lotka--Volterra model}\index{Lotka--Volterra
  model}.  The Lotka--Volterra model,
\begin{align*}
R' &= aR-bRF, \\
F' &= cRF-dF,
\end{align*}
while it had a major impact on the development of mathematical
biology, was found to be flawed in several important ways.  The chief
problem is that the equilibrium point $(d/c, a/b)$ is a neutral
equilibrium point---given any starting state, the system would follow
a closed trajectory.  This in itself was all right and, in fact,
stimulated a great many important investigations on whether or not
cycles were an intrinsic feature of many populations.  The difficulty
was that there were so many possible closed trajectories---which one
the system followed depended on where it started.  A second
difficulty, related to the first, is that there is a first integral
for the Lotka--Volterra model.  What is seen as a virtue in a physical
system like the pendulum---since it is equivalent to the conservation
of energy---is unrealistic in an ecological system, where there are
almost certainly too many outside forces at work for any quantity to
be conserved there.  In the following exercises we will explore some
of these behaviors.  As before, you can either perform a general
analysis of the model or use the specific parameter values $a=.1$,
$b=.005$, $c=.00004$, and $d=.04$.  

\sub Sketch the vector field, together with some typical trajectories,
in the rest of the $R$-$F$ plane, including negative values. What
happens to any trajectory starting at a state with a negative $R$ or
negative $F$ value?  

\sub For this exercise you will need to go back to a computer program
that implements Euler's method of approximating the trajectory by
drawing a straight line segment from a point in the direction
indicated by the velocity vector (commercial packages use fancier
routines which accommodate for the kind of phenomena you are about to
see!). Using the specific values for $a$, $b$, $c$, and $d$ suggested
above, starting from the point $(2000, 10)$ in the $R$-$F$ plane, and
using a time step $\Delta t=1$, draw the first 500 segments of Euler's
approximation to the trajectory.  What does the trajectory look like?
Would you think the trajectory was a closed loop on the basis of this
result?  How small does $\Delta t$ have to be before the trajectory
looks like it closes?  Can you explain this phenomenon?  \sub Using
the same values for $a$, $b$, $c$, and $d$ as in the preceding part,
start at the point $(2000, 1)$ and use $\Delta t=2$.  This time
calculate the first 1000 segments of Euler's approximation; what
happens? (Your computer will probably give you some sort of overflow
message.) Can you explain this? (Think about your answer to part (a).)

\sub {\bf Getting a first integral for the system}\quad Show
that the Lotka--Volterra equations imply that
\[
\dfrac{R'}{R}(cR-d) = \dfrac{F'}{F}(a-bF).
\]
Integrate this equation and show that the expression
\[
cR+bF-d\ln{R} - a \ln{F}
\]
must be a constant for all points on a given trajectory.  If we know
one point on the trajectory (such as the starting point), we can
evaluate the constant.  \sub Show that the function $f(R)=cR-d\ln R$
is decreasing for $0<R<d/c$ and is increasing for
$d/c<R<\infty$. Hence argue that for any given value of $F$ there are
at most two values of $R$ giving the same value for the expression
$cR+bF-d\ln R-a\ln F$.  Hence conclude that the trajectories for the
Lotka--Volterra equations can't be spirals, but must then be closed
loops.

\subsubsection{The pendulum}

\vspace{-10pt}

\num Suppose instead of an idealized frictionless pendulum, we wanted
to model a pendulum that ``ran down''.  One approach we might try is
to throw in a term for air resistance.  Let's see what happens when we
add a term to the expression for $v'$ which suggests that there is a
drag effect which is proportional to the value of $v$---the larger $v$
is, the greater will be the drag.  Here are equations that do this:
\begin{align*}
x' &= v, \\
v' &= -\sin x-.1v.
\end{align*}
Perform a vector field analysis of this model, indicating the regions
where the velocity vectors are pointing in the various combinations of
up, down, right, and left.  Try sketching in some trajectories.  Where
are the equilibrium points? What kinds are they?



\subsubsection{The Anderson--May model}

\vspace{-10pt}

\num Consider $dE_1/dt=\Lambda_1-\mu_1E_1+a_1E_1E_2/(1+b_1E_1E_2)$.
For what values of $E_1$ is it possible to find a value for $E_2$
making $dE_1/dt=0$?  Express your answer in terms of the parameters
$\Lambda_1$, $\mu_1$, $a_1$, and $b_1$.  Is your answer consistent
with the graph on page~\pageref{`amprob'}?

\num In\label{AMay2} the same book---{\em Infectious Diseases of
  Humans}---containing the previous model, Anderson and May propose
another model to explain the acquisition of (apparently) permanent
immunity.  In this model there is just the virus and the lymphocyte
cells (effector cells) that kill the virus.  We denote their
populations at time $t$ by $V(t)$ and $E(t)$.  They propose the model
\begin{align*}
dE/dt &= \Lambda-\mu E +\varepsilon VE, \\
dV/dt &= rV-\sigma VE.
\end{align*}
Here $\Lambda$ is the (constant) rate of background production of the
lymphocytes by the bone marrow, $\mu$ is the per capita death rate of
such cells, and $r$ is the intrinsic growth rate of the virus if none
of the specific lymphocytes was present.  Both the increased
production of the lymphocytes and the death of the virus are assumed
to proceed at rates proportional to the number of their interactions,
determined by their product.

\sub Show that in the absence of any virus, the effector cells have a
stable equilibrium of $\Lambda /\mu$.

\sub Perform a state space analysis of the vector field.  Note that
there will be two very different cases, depending on whether $\Lambda
/\mu>r/\sigma$ or $\Lambda /\mu<r/\sigma$.  In each case say what you
can about the equilibrium points and the expected long-term behavior
of the system.

\sub Using parameter values $\Lambda=1$, $\mu=r=.5$, and
$\varepsilon=\sigma=.01$, and starting values $E=V=1$ find the
resulting trajectory. (The trajectory will be a spiral, but it moves
in very slowly.)

\sub How long, approximately, will it take the spiral to make one
revolution? If this time, call it $T$, is roughly the same length as
the lifetime of the infected individual, what will appear to be
happening?  It might help to plot both $E$ and $V$ as functions of
time over the interval $[0,T]$.


%%%%%%%%%%%%%%%%%%%%%%%%%%%%%%%%%%%%%%%%%%%%%%%%%%%%%%%%%%%%%%%%%%%%%%%%%
%                                                                       %
%                                Section 2                              % 
%                                                                       %
%%%%%%%%%%%%%%%%%%%%%%%%%%%%%%%%%%%%%%%%%%%%%%%%%%%%%%%%%%%%%%%%%%%%%%%%%

\section{Local Behavior of Dynamical Systems}
\label{sec-eqs}

\subsection{A Microscopic View}

One of the themes of this book has been the concept of the ``microscope''.  
%
\mar{Phase portraits under the microscope}
%
When we zoom in on some part of a geometrical object, the structure
typically becomes much simpler. In chapter 3 we used this approach to
think about the behavior of functions.  In this section we will use
the same idea to analyze the behavior of a vector field and its phase
portrait. There are two parts to this process:
\begin{enumerate}

\item We shift the origin of the coordinate system to center on the
  point we are interested in---we {\bf localize}\index{localize}---and

\item We approximate both the vector field and its phase portrait by
  suitable linear approximations---we {\bf
    linearize}\index{linearize}.

\end{enumerate}
 
To get a feel for how this works, let's go back and look at problem 4
on page~\pageref{`basins'} of the previous section.  There we had two
species X and Y competing for the same food source.  We modeled the
dynamics of this system by the equations\label{`saddle'}
\begin{align*}
x' &= .15(1 - .005x - .010y)x, \\
y' &= .03(1 -.006x - .005y)y.
\end{align*}
The phase portrait for this system looks like the following figure.

\vspace{225pt}

\special{psfile=../ps/calc2/locbeha1.ps}

The three equilibrium points---$P=(0,200)$, $R=(1000/7, 200/7)$, and
$S=(0, 200)$---are indicated, together with a generic point $Q=(35,
50)$.  Note that $P$ and $S$ are attractors and that $R$ is a saddle
point.  As was the case with the Anderson--May model, there is a
trajectory flowing away from $R$ to each of the attractors.  There are
also two trajectories (not shown) flowing directly toward $R$ and
forming the boundary between the basins of attraction for $P$ and $S$.
We will see how to construct this boundary shortly
(page~\pageref{`boundary'}).

Let's first zoom\label{`andmay'} in on the point $Q$ and see what the
phase portrait looks like there.  If we take the region $\pm 1$ unit
on either side of $Q$, we get the following phase portrait:

\vspace{110pt}

\special{psfile=../ps/calc2/locbeha2.ps}

At this level, all the trajectories appear to be parallel straight
lines. How could we have anticipated this picture?  The first step in
analyzing this phase portrait is to observe that since we are
interested in its behavior near $Q=(35, 50)$, instead of working with
the variables $x(t)$ and $y(t)$, we introduce new variables $r(t)$ and
$s(t)$ which measure how far we are from $Q$:
\begin{align*}
r(t) &= x(t)-35, \\
s(t) &= y(t) - 50.
\end{align*}
The effect of this transformation is simply to shift the origin to the
point $Q$---the location of every point in the plane is now measured
relative to $Q$
%
\mar{Shifting the origin}
%
rather than to the $x$-$y$ origin. A point is close to the point $Q$
if its $r$-$s$ coordinates are small. Further, if we are given the
$r$-$s$ coordinates of a point, we can always recover the $x$-$y$
coordinates, and vice versa---we can transform in either direction:
\begin{alignat*}{2}
r &= x-35, &\quad\Longleftrightarrow\quad x &= r+35, \\
s &= y-50, &\quad\Longleftrightarrow\quad y &= s+50.
\end{alignat*}

Next, note that $r'(t) = x'(t)$ and $s'(t) = y'(t)$ so that the new
variables change at the same rates as the old ones.  We can now
express our original differential equations in terms of the variables
$r$ and $s$ by replacing $x'$ by $r'$, $x$ by $r+35$, $y'$ by $s'$,
and $y$ by $s+50$.  When we do this, we get
\begin{align*}
r' & =  .15(1- .005(r+35) - .010(s+50))(r+ 35)\\
   &= 1.70625+.0225r-.0525s-.00075r^2-.0015rs, \\
s' &= .03(1-.006(r+35) - .005(s+50))(s+50)\\
   &= .81-.009r+.0087s-.00018rs-.00015s^2.
\end{align*}

What we have accomplished by this is to transform a problem about
trajectories near the point $(35, 50) )$ in the $x$-$y$ plane into a
problem about trajectories near the origin in the $r$-$s$ plane---we
have {\bf localized} the problem to the point we are interested in.

The second step comes in analyzing the $r$-$s$ system: since we are
only interested in its behavior near the origin, we will be looking at
values of $r$ and $s$ that are small.  Under these
%
\mar{Near an ordinary point, a vector field is\\ almost constant}
%
circumstances, the contributions of the constant terms will far
outweigh the contributions of any of the terms involving $r$ and $s$.
For instance, in our current example we are looking at a window that
is $\pm 1$ unit wide and $\pm 1$ unit high around $Q$. In this window,
the terms involving $r$ or $s$ are at most 3\% of the constant term in
the case of $r'$, and a little over 1\% in the case of $s'$. If we had
used a smaller window, the contributions of the non-constant terms
would be even less significant.  This means that {\em near the $r$-$s$
  origin\/} the vector field for this system is well-approximated by
the behavior of the related {\bf constant linear system}:
\begin{align*}
r' &= 1.70625, \\
s' &= 0.81.
\end{align*}
Note that 1.70625 and .81 are just the values of $x'$ and $y'$ at $Q$. 

In this linearized system, any change $\Delta t$ in the time produces
a change $\Delta r =1.70625\Delta t$ in $r$, and a change $\Delta
s=.81\Delta t$ in $s$. Thus the velocity vectors in the vector field
near $Q$ would all have the same length
%
\mar{Near an ordinary point all trajectories look\\ the same}
%
and would be pointing in the same direction, with slope $\Delta
s/\Delta r =.81/1.70625=.4747$. This in turn means that near $Q$
all trajectories have the same slope and are traversed at the same
speed.
        
We would see a similar picture---a family of parallel straight
lines---whenever we zoom in on the phase portrait near any other
ordinary (i.e., non-equilibrium) point $(x_\ast, y_\ast)$ The vector
field near such a point can always be approximated by a constant
linear system of the form
\begin{align*}
r' &= e, \\
s' &= f,
\end{align*}
where $e$ and $f$ are the values of $x'$ and $y'$ at $(x_\ast,
y_\ast)$.  The trajectories of this approximating linear system will
be lines of slope $f/e$.

Near an 
%
\mar{Equilibrium points\\ are different}
%
equilibrium point, the picture is more complicated.  No matter how far
in we zoom, the phase portrait never looks like a family of straight
lines.  For instance, here's what the picture looks like when we zoom
in on $R=(1000/7, 200/7)\approx (142.857, 28.571)$\,:

\vspace*{120pt}
\special{psfile=../ps/calc2/locbeha3.ps}

\noindent
---if we zoomed in to a window 1/100-th the size of this one, the
picture would be indistinguishable from this one.

Here we see four trajectories that look almost like straight
lines---two coming directly towards $R$ and two going directly away.
All the other trajectories appear to be asymptotic to these two sets.
On page~\pageref{fixedline} in the next section you will see how to
find the equations of these asymptotes.

What happens when we linearize the vector field at $R$?  As before, we
first shift the origin so that it is centered at $R$ by changing to
coordinates $r$ and $s$, where
\begin{align*}
r(t) &= x(t)-1000/7, \\
s(t) &= y(t) - 200/7.
\end{align*}

When we then write the differential equations in terms of $r$ and $s$,
we get as before that $x'=r'$ and $y'=s'$ and
\begin{align*}
r' &= .15(1- .005(r+1000/7) - .010(s+200/7))(r+ 1000/7)\\
   &= -.107143r-.214286s-.00075r^2-.0015rs,  \\
s' &= .03(1-.006(r+1000/7) - .005(s+200/7))(s+200/7)\\
   &= -.00514286r-.00428571s-.00018rs-.00015s^2.
\end{align*}

This time, though, the constant term 
%
\mar{Linear approximation of the vector field}
%
in the expression for both $r'$ and $s'$ is 0.  This is because the
point $R$ was an equilibrium point, which meant that both $x'$ and
$y'$, and hence $r'$ and $s'$, were 0 there.  If we are considering
only small values of $r$ and $s$, though, say much smaller than 1,
then the terms involving $r^2$ or $s^2$ or $rs$ will be much smaller
than the terms involving $r$ and $s$ alone.  We can therefore simplify
our equations at $R$ by taking only the first powers of $r$ and $s$,
getting for the linearized system
\begin{align*}
r' &=  -.107143r-.214286s,  \\
s' &=  -.00514286r-.00428571s.
\end{align*}

In a similar fashion we could hope to explore the behavior of any
other dynamical system about any of its equilibrium points by
approximating the vector field there by a linear system of the form
\begin{align*}
r' &= ar +bs, \\
s' &= cr + ds,
\end{align*}
for suitable constants $a$, $b$, $c$, and $d$. 

We will see in section~\ref{taxonomy} how to use this linearized form of the vector field to discover many of the properties of equilibrium points.

How can we find values for the constants $a$, $b$, $c$, and $d$?  If
the differential equations specifying the rates of change of the
variables are polynomials, then we can proceed as above:
\begin{itemize}

\item Shift the origin to the point we're interested in;

\item Express the rate equations in terms of the new local variables;

\item Throw away all the terms except the first degree terms.

\end{itemize}

This process requires some fairly tedious algebra.  Moreover, what if
the differential equations are not polynomials?  Suppose, for
instance, we wanted to study the local behavior of the Anderson--May
model (page~\pageref{`amprob'}) at the saddle point $P_2=(92.0869,
92.0869)$. Note that the differential equations are of the form
\begin{align*}
dE_1/dt &= f_1(E_1,E_2), \\
dE_2/dt &= f_2(E_1,E_2),
\end{align*}
where $f_1$ and $f_2$ are the functions given in the text. But $f_1$
and $f_2$ are just
%
\mar{To linearize a vector field, linearize the functions that determine it}
%
functions, and we learned in chapter 3 how to construct locally linear
approximations to them.  This was, in fact, how we defined derivatives
in the first place.  Thus if $E_1$ changes by a small amount $\Delta
E_1=E_1-92.0869$, the function $f_i$ will change by approximately
$\partial f_i/\partial E_1 \times \Delta E_1$.  Similarly, a small
change $\Delta E_2=E_2-92.0869$ will produce a change of approximately
$\partial f_i/\partial E_2 \times \Delta E_2$ in the function $f_i$.
The total change in the function $f_i$ can then be approximated by the
sum of these changes:
\begin{align*}
\Delta f_1(E_1,E_2) &\approx 
       \dfrac{\partial f_1}{\partial E_1}\Delta E_1 
       + \dfrac{\partial f_1}{\partial E_2}\Delta E_2, \\
\Delta f_2(E_1,E_2) &\approx 
       \dfrac{\partial f_2}{\partial E_1}\Delta E_1 
       + \dfrac{\partial f_2}{\partial E_2}\Delta E_2.
\end{align*}

But since $P_2$ is an equilibrium point, 
%
\mar{General form for the local linearization at\\ an equilibrium point\ldots}
%
we have by definition that $f_1$ and $f_2$ are both zero there, so
$\Delta f_i(E_1,E_2) = f_i(E_1,E_2) - f_i(P_2)$ is just
$f_i(E_1,E_2)$.  Further, if you look closely you will see that the
quantity $\Delta E_1=E_1-92.0869$ is identical with what we have been
calling the local coordinate $r$, and $\Delta E_2=E_2-92.0869$ is just
the other local coordinate $s$.  Thus, since $E_1'=r'$ and $E_2'=s'$,
we have
\begin{align*}
r'&= \dfrac{\partial f_1}{\partial E_1}\,r 
     + \dfrac{\partial f_1}{\partial E_2}\,s, \\
s'&= \dfrac{\partial f_2}{\partial E_1}\,r 
     + \dfrac{\partial f_2}{\partial E_2}\,s,
\end{align*}
where the partial derivatives are evaluated at $P_2$.  Notice that
there is nothing in this expression which is specific to this
particular problem.  The local linearization of any vector field at
any equilibrium point will be in this form.

Finally, using the values given for the different parameters back on
page~\pageref{`amprob'}, we can evaluate all the partial derivatives
to get the specific local linearization for the point $P_2$:
\begin{align*}
r' &= -94.5525\,r + 905.448\,s, \\
s' &= 905.448\,r - 94.5525\,s,
\end{align*}

We will see in the next section how knowing this form will allow us to
find the boundary between the two basins of attraction.

For completeness, 
%
\mar{\ldots and at\\ a generic point}
%
let's remind ourselves of what the local linearization would look like
at a nonequilibrium point in the current formulation.  The result is
immediate and simple, using the analysis we used before.  If $Q$ is a
generic point, then the local linearization consists of parallel
lines, whose slopes are given by the constant rate equations
\begin{align*}
r' &= f_1(Q), \\
s' &= f_2(Q).
\end{align*}

\newpage

\subsection{Exercises}

\num Find the local linearizations at all the equilibrium points in
exercises 2--4 at the end of the previous section.

\numa  The Lotka--Volterra equations
\begin{align*}
R' &= aR-bRF, \\
F' &= cRF-dF,
\end{align*}
have an equilibrium point at $(R, F) = 00(d/c, a/b)$.  

\sub What is the local linearization there?  

\sub What is a striking feature of this linearization, and what is its
physical significance?

\sub The trajectories for the local linearizations turn out to be
ellipses.  If $r$ and $f$ are the local variables, find constants
$\alpha$ and $\beta$ such that the expression $\alpha\,
r^2+\beta\,f^2$ is constant on any trajectory.

\num Find the local linearization at the point $P_1$ in the
Anderson--May model for the acquisition of immunity discussed in the
previous section, using the parameter values given in the text on
page~\pageref{paramvals}.

\num Go back to the second Anderson--May model analyzed in problem~8
of the previous section (page~\pageref{AMay2}).  Using the parameter
values given in part~(c) there, find the local linearizations at all
equilibrium points.

%%%%%%%%%%%%%%%%%%%%%%%%%%%%%%%%%%%%%%%%%%%%%%%%%%%%%%%%%%%%%%%%%%%%%%%%%
%                                                                       %
%                                Section 3                              % 
%                                                                       %
%%%%%%%%%%%%%%%%%%%%%%%%%%%%%%%%%%%%%%%%%%%%%%%%%%%%%%%%%%%%%%%%%%%%%%%%%

\section{A Taxonomy of Equilibrium Points} 
\label{taxonomy}

In the exercises and examples we have seen so far in this chapter, 
%
\mar{An intuitive classification of equilibrium points} 
%
there have been several kinds of trajectories near equilibrium points:
spirals towards and spirals away from the equilibrium, closed loops
about the equilibrium, trajectories that looked vaguely like
hyperbolas, and trajectories that seemed to arc more or less directly
into or away from the equilibrium.  It turns out that this rough
classification covers virtually all the equilibrium behaviors we might
encounter in a two-dimensional state space.  There are many ways to
demonstrate this, but we can accomplish almost everything with a
couple of simple insights.  We begin with a summary of the different
kinds of equilibrium points, then turn to the question of devising
ways to figure out from the equations what kind we are dealing with.

Suppose, then, that we are studying a two-dimensional dynamical system
and that we have linearized the system at an equilibrium point.  The
point is either an attractor, repellor, saddle point, or neutral
point.  Attractors and repellors can be further subdivided according
to whether they have one or two straight line trajectories, or whether
their trajectories are spirals.  Note that any attractor can be
converted into a repellor simply by reversing the arrows, and vice
versa (how do you accomplish this arrow reversal at the level of the
defining differential equations?).  If you reverse all the arrows at a
saddle point, you get another saddle point.  If you reverse the arrows
at a neutral point, you get the same closed loops, but they are
traversed in the opposite direction.
        
Here, then, is a listing of all the kinds of equilibrium points. There
are five generic types.  ({\em Generic\/} here means ``general''; if
you generate a random equilibrium point, it will almost certainly be
one of these.)  They are most easily categorized by whether or not
they have {\bf fixed line}\index{fixed line} trajectories---that is,
trajectories which are straight lines going directly toward or
directly away from the equilibrium point.

\aside{The existence or not of straight line trajectories and how to
  find them when they do exist is an instance of the so-called {\em
    eigenvector\/} problem.  Analogous problems occur elsewhere in
  many parts of mathematics, physics, and even population biology.
  Being able to find such eigenvectors efficiently is an important
  problem in computational mathematics.}

\smallskip      
\noindent
{\bf Nodes}. \index{node}\index{equilibrium point!node} Two pairs of
fixed lines, all trajectories flowing toward the equilibrium
(attractors) or away from it (repellors).

\vspace*{1.8in}

\special{psfile=../ps/calc2/classeq1.ps}

\smallskip
\noindent
{\bf Spirals}. \index{spiral}\index{equilibrium point!spiral} No
fixed lines, all trajectories spiraling toward the equilibrium
(attractors) or away from it (repellors).

\vspace*{1.8in}
\special{psfile=../ps/calc2/classeq2.ps}

\smallskip
\noindent
{\bf Saddle Point}. \index{saddle point}\index{equilibrium point!saddle point} Two pairs of fixed lines, with the flow along one pair
being toward the equilibrium, and the flow along the other pair away
from it.  All other trajectories are asymptotic to these lines.

\vspace*{1.8in}
\special{psfile=../ps/calc2/classeq3.ps}

In addition to these five generic cases, there are three more types
which arise under more specialized conditions:

\smallskip
\noindent
{\bf Special Nodes}.  One pair of fixed lines, all trajectories flowing toward the equilibrium (attractors) or away from it (repellors). 

\vspace*{1.8in}
\special{psfile=../ps/calc2/classeq4.ps}

\smallskip
\noindent
{\bf Center}. No fixed lines, all trajectories flowing around the
equilibrium in closed loops.  \index{neutral point}\index{equilibrium
  point!neutral}

\vspace*{1.8in}
\special{psfile=../ps/calc2/classeq5.ps}

Except for a variety of highly specialized (or {\em degenerate}, in
mathematical terminology) cases, examples of which are given in the
exercises, the region near every equilibrium point will look like one
of the above (although the exact shape may vary).

Clearly, it would be helpful to have an efficient way to determine
whether or not fixed lines exist, and what their equations are if they
do.

\subsection{Straight-Line Trajectories}

Given a dynamical system
\begin{align*}
r' &= ar +bs,  \\
s' &= cr + ds,
\end{align*}
how can we tell whether or not it has any straight-line
trajectories?\label{fixedline} If $b = 0$, then the (vertical) line $r
= 0$ is a trajectory.  Otherwise, note that the line $s=mr$ will be a
trajectory for this system provided the slope of the line---namely
$m$---equals the slope of the vector field at every point $(r,s)$ on
the line.  But the slope of the vector field at any point $(r,s)$ is
just $s'/r'$, which in turn is equal to $(cr + ds)/(ar +bs)$.  Since
every point on the line of slope $m$ is of the form $(r, mr)$, what we
are really asking, then, is whether there are any
%
\mar{\vspace{20pt}The condition for a fixed-line trajectory}
%
values of $m$ which satisfy the equation
        $$
        m=\frac{cr+dmr}{ar+bmr}=\frac{c+dm}{a+bm}
        $$

To see how this works, let's return to the example of two competing
species which we last looked at on page~\pageref{`saddle'}.  There we
zoomed in on the saddle point $R=(1000/7, 200/7)$ and found that the
local linear approximation was
\begin{align*}
r' &= -.1071r -.2143s, \\
s' &= -.0051r - .0043s.
\end{align*}
If this system has a straight-line trajectory of slope $m$, then $m$
must satisfy
        $$      
        m= \dfrac{ -.0051 - .0043m}{ -.1071 -.2143m},
        $$
which leads to the quadratic equation
        $$
        .2143m^2+.1028m-.0051=0,
        $$
which has roots
        $$
        m=.0454 \quad \mbox{and} \quad m = -.5250.
        $$ 
Thus the lines $s=.0454r$ and $s=-.5250r$ are trajectories of the
linear system.  To be more exact, each of these lines is made up of
three distinct trajectories: the portion of the line consisting of all
points with $r>0$, the portion with $r<0$, and the origin (which is
the saddle point $R$) by itself, which is always a trajectory in any
linear system.  To see whether flow along these trajectories is
towards the origin or away from it, we could look to see where the
lines lie in the state plane.  It is just as simple, though, to try a
test point.  For instance, a typical point on the line $s=.0454r$ is
$(1, .0454)$.  When we substitute these values into the original rate
equations, we find that
\begin{align*}
r' &= -.1071\times 1-.2143\times .0454, \\
s' &= -.0051\times 1 - .0043 \times .0454.
\end{align*}

We don't even need to do the arithmetic to be able to tell that both
$r'$ and $s'$ are negative at this point, hence both $r$ and $s$ are
decreasing, which means that on the line of slope .0454 movement is
towards the origin.  Similarly, on the line of slope $-.5250$ the flow
is away from the origin.  Finally, it turns out (as is the case with
every linear system with straight-line trajectories) that every other
trajectory is asymptotic to these lines.

The crux of this approach was the use of the quadratic formula.  Of
course, it may happen---and we will see examples in the
exercises---that when we try the same approach on another system we
find there are no real roots to the equation.  This means that there
are no fixed lines, so that trajectories must be spirals or closed
loops.

\subsubsection{Attractors and Basins of Attraction}\label{`boundary'}

One byproduct of the analysis in the previous section is that it gives
us a technique for sketching the boundary separating two basins of
attraction. Let's continue with the previous example to illustrate how
this is done. We observed that the boundary between the two basins was
formed by the two trajectories coming directly into the saddle point
$R$ between the two attractors $P$ and $S$.  We have just seen that
near $R$ these two trajectories looked like the straight line of slope
.0454.  We can therefore take a point on this line on each side of $R$
and run the system backward (if we go forward, we simply approach $R$)
in time to reconstruct the trajectories, and hence get the boundary of
the basins of attraction.


\subsection{Exercises}


\num In this exercise we look at a number of different linear systems
to see what kinds of trajectories we get.  In each case you should
sketch the trajectories.  Do this as before by first identifying the
regions in the plane where $r'=0, r'>0$, and $r'<0$, and similarly for
$s'$.  Then sketch trajectories consistent with this information.  You
might want to use a graphing program to check any answer you're unsure
of.

\sub $r' = 4r + s$, \hspace{.25in} $s' = 2r + 3s$.

\sub $r' = 4r + s$, \hspace{.25in} $s' = -2r + 3s$.

\sub $r' = 2r + 3s$, \hspace{.25in} $s' = 4r + s$.

\sub $r' = -4r + 4s$, \hspace{.25in} $s' = 2r + s$.

\sub $r' = -.4r - 4s$, \hspace{.25in} $s' = 2r - .5s$.

\sub $r' = -.4r - 4s$, \hspace{.25in} $s' = 2r + .4s$.

\sub   Make up and analyze four more linear systems.

\num If you start with a given linear system and consider the related
system in which all the coefficients are four times as big, how do the
trajectories change?

\num If you start with a given linear system and consider the related
system in which all the coefficients have their signs reversed, how do
the trajectories change?


\numa Use the quadratic formula to find the general solution to the
equation
        $$
        m=\frac{c+dm}{a+bm}.
        $$

\sub In exercises 1, 3, and 4 in the previous section you found local
linearizations at the equilibrium points of a number of examples
discussed earlier.  Determine which of these have straight-line
trajectories and which do not.  For those that do, find the equations
of the lines and determine for each line whether the flow is towards
the origin or away from it.

\sub What is the general condition for a linear dynamical system to
have straight-line trajectories?

\num Make up a system that has the lines of slope $\pm 1$ as
trajectories.

\num What is the condition for a system to have exactly one fixed
line?  Construct a couple of systems that have only one fixed line and
sketch their phase portraits.

\num {\bf Degeneracy}\quad The analysis developed in this section
implicitly assumed that in the local linearization, at least one of
the coefficients in each of the expressions for $r'$ and $s'$ was
non-zero.  If this is not true, then many more possibilities open up.
The following two systems have the origin as their only equilibrium
point.  In each case, write down the local linearization and draw in
the trajectory pattern for the linearized system.  Notice that the
linearized systems have more than one equilibrium point.  Then do the
standard phase plane analysis for the original system---identify the
regions in the plane where $r'=0$and where $s'=0$, and specify what
the direction field is doing in the rest of the plane, as usual.
Sketch in some typical trajectories.  Comment on the connections
between the linearized and unlinearized forms.  

\sub $r' = r^2$, \hspace{.25in} $s' = -s$.  You should see sort of a
hybrid between a saddle point and an attractor here.

\sub $r' = r^2+s^2$, \hspace{.25in} $s' = r$.


\numa Use the technique presented at the end of this section
(page~\pageref{`boundary'} to graph the boundary between the two
basins of attraction.

\sub In the same way, construct the boundary between the two basins in
the competing species model we've been discussing---problem 4 on
page~\pageref{`basins'}.

\subsubsection{Distance from the Origin}

Another way to distinguish between different kinds of trajectories is
to see how their distance from the origin varies over time. For saddle
points the distance will first decrease and then increase.  For spiral
attractors and nodal attractors the distance may be always decreasing,
or it may fluctuate, depending on how flat the trajectory is.

Again, let's look at a general linear system
\begin{align*}
r' &=  ar +bs,  \\
s' &=  cr + ds.
\end{align*}
onsider the system moving along some trajectory in $r$-$s$ space.  At
time $t$ it will be at a point $(r(t), s(t)$, situated at a distance
$d(t) = \sqrt{r(t)^2 + s(t)^2}$.  We would like to know how the
function $d(t)$ behaves.  Is it always increasing?  Always decreasing?
Or does it have local maxima and minima?  To answer this we need to
know if $d'(t)$ is ever $=0$, or if it is always positive or always
negative.  We can simplify our calculations if we look at the square
of the distance: $D(t) = d(t)^2 =r(t)^2 + s(t)^2$.  The function $D$
will be increasing and decreasing at exactly the same points as the
function $d$, and it's easier to work with.

\numa Show that
\begin{align*}
D'(t) &= 2 r(t)r'(t) + 2 s(t)s'(t)\\
      &= 2[r(ar+bs)+s(cr+ds)]\\
      &= 2[ar^2+(b+c)rs+ds^2].
\end{align*}

\sub Show that if we look at points on the line of slope $m$, so that
$s=mr$, we will have $D'(t)=0$ there if and only if
        $$
        a+(b+c)m+dm^2=0\, .
        $$

\sub Use the quadratic formula to conclude that this happens precisely
where
        $$
        m=\frac{-(b+c) \pm \sqrt{(b+c)^2-4ad}}{2d}\, .
        $$

\sub Show in particular, if $(b+c)^2-4ad<0$, there are no solutions to
$D'(t)=0$, and the distance must always be strictly increasing along
all trajectories, or strictly decreasing along all trajectories.

\sub    Return to the example
\begin{align*}
r' &=  -.1071r -.2143s, \\
s' &=  -.0051r - .0043s,
\end{align*}
and find the equations of the two lines where trajectories pass
closest to the origin.  These lines will not be trajectories
themselves.  Their significance is that the `vertices' of all the
trajectories will lie along them.

\num Choose four of the exercise in the first part of this section and
analyze them to see where (and whether) trajectories have a closest
approach to the origin.

\numa Use the results of this section to construct a dynamical system
whose trajectories are spirals that are always moving away from the
origin.

\sub Use the results to construct a dynamical system whose
trajectories are flattened spirals, so that the distance from the
origin, while increasing overall, has local maxima and minima.

\num It turns out that trajectories which form closed loops should
really be considered as a special kind of spiral.  In fact, a
flattened spiral will close up precisely when the two directions in
which the distance is a maximum or minimum are perpendicular to each
other.  Express this as a condition on the coefficients $a, b, c, $
and $d$ in the dynamical system.

\num Write down the equations of some dynamical systems that will have
closed orbits.

%%%%%%%%%%%%%%%%%%%%%%%%%%%%%%%%%%%%%%%%%%%%%%%%%%%%%%%%%%%%%%%%%%%%%%%%%
%                                                                       %
%                                Section 4                              % 
%                                                                       %
%%%%%%%%%%%%%%%%%%%%%%%%%%%%%%%%%%%%%%%%%%%%%%%%%%%%%%%%%%%%%%%%%%%%%%%%%

\section{Limit Cycles}
\label{sec-repel}

\enlargethispage*{20pt}
With this analysis of the behavior of vector fields near equilibrium
points, we now know most of the possibilities for the long-term
behavior of trajectories.  The one important phenomenon we haven't
discussed is {\bf limit cycles}.\index{limit cycle} To see an example
of this, let's return to May's predator--prey model we first
encountered in chapter 4.  If $x(t)$ and $y(t)$ are the prey and
predator populations, respectively, at time $t$, then the general form
of May's model is
\begin{align*}
x' &= ax\left(1-\frac{x}{b}\right) - \frac{cxy}{x+d},  \\
y' &= ey\left(1-\frac{y}{fx}\right);
\end{align*}
the parameters $a$, $b$, $c$, $d$, $e$ and $f$ are all positive.
 
Using parameter values of $a = .6$, $b = 10$, $c = .5$, $d = 1$, $e =
.1$ and $f = 2$, let's take several different starting values and
sketch the resulting trajectories.  Here's what we find:

\vspace*{2.5in}
\special{psfile=../ps/calc2/limcyc.ps}

Notice that no matter where we start, the trajectory is apparently
always drawn to the closed loop shown in dashes above.  This loop
is an example of an {\bf attracting limit cycle}.  As usual, we could
reverse all the arrows in our vector field, in which case this example
would be converted to a {\bf repelling limit cycle}.

A limit cycle is very different from the kind of behavior we saw in
the neighborhood of a neutral equilibrium point called a center.
Around a center there is a closed loop 
%
\mar{Limit cycles give models for cyclic behavior} 
%
trajectory through every point: displace the state slightly, and it
would move happily along the new loop.  If the state is on an
attracting limit cycle, though, and you displace it, it will move back
toward the cycle it started from.  For this reason limit cycles make
very good models for cyclic behavior, whether it is in the firing of
neurons or population cycles of mammals.

The size of the limit cycle, and even its very existence, depends on
the specific values of the parameters in the model.  If you change the
parameters, you change the limit cycle.  If you change the parameters
enough, the limit cycle may disappear all together. (See the
exercises.)

A result proved early in this century is the {\bf
  Poincar\'{e}--Bendixson Theorem} \index{Poincar\'{e}--Bendixson
  Theorem} which says that equilibrium points and limit cycles are as
complicated as dynamical systems in two variables can get.  Once we
pass to three variables, the situation becomes much more complicated.
Many of the phenomena associated with such systems have been
discovered only within the past 50 years, and their exploration is a
subject of continuing research.  In the next section we will give a
brief introduction of some of the new behaviors that can arise.


\subsection{Exercises}

\num Using the parameter values given in the text, find the
coordinates of the equilibrium point at the center of the limit cycle
and show that it is, in fact, a repellor.

\num May's model is interesting in that it exhibits a phenomenon known
as {\bf Hopf bifurcation}.\index{Hopf bifurcation} Namely, the
existence of a limit cycle depends on the values of the parameters.
Choose one of the parameters in May's model and try a range of values
both larger and smaller than in the example we've worked out. At what
value does the limit cycle disappear?  When this happens, the
equilibrium point inside the cycle has become an attractor.  Can you
work out analytically when this happens?

%%%%%%%%%%%%%%%%%%%%%%%%%%%%%%%%%%%%%%%%%%%%%%%%%%%%%%%%%%%%%%%%%%%%%%%%%
%                                                                       %
%                                Section 5                              % 
%                                                                       %
%%%%%%%%%%%%%%%%%%%%%%%%%%%%%%%%%%%%%%%%%%%%%%%%%%%%%%%%%%%%%%%%%%%%%%%%%

\section{Beyond the Plane:\\ Three-Dimensional Systems}

Up to now we have worked with dynamical systems in which there are
only two interacting quantities.  We have thought of the two
quantities as specifying a point in the state space, which we think of
as some subset of the plane. The dynamical system defined a vector
field on the state space.  These geometric notions carry over to
dynamical systems involving more than two interacting quantities.

In particular, if we have a dynamical system consisting of three
interacting quantities, then we think of the values of the three
quantities as specifying a point (or state) in three-dimensional
space.  So, for instance, if we have an ecological system consisting
of three species, then we think of the numbers $x$, $y$, $z$ of each
of the three species as specifying a point $(x,y,z)$ in space.  The
set of all possible points or ``states'' is the set of points
\[
\{(x,y,z): x \ge 0, y\ge 0, z\ge 0\}
\]
that constitute the ``first octant'' in Cartesian 3-space.  We think
of the dynamical system as a vector field: that is, as a rule which
assigns to each point of the state space a vector.  As in the case of
the plane, we can define equilibrium points, trajectories, limit
cycles, attractors, and the like.

In three-space there is a much wider range of behavior possible, even
in the case of equilibrium points.  We do, of course, have point
attractors and repellors: all trajectories near a point attractor flow
towards the attractor and all trajectories near a point repellor flow
away from the repellor.  However, a greater range of combinations is
possible: an equilibrium point can attract all points along some
plane, but repel all other points.  Or the equilibrium point could be
a center, surrounded by closed orbits lying in some plane which
attract trajectories off the plane.
  
\vspace*{2.4in}
\special{psfile=../ps/calc2/ctr.ps}

\enlargethispage*{15pt}
It is worth pointing out that we can represent the two-dimensional
systems we've been exploring so far in this chapter in three
dimensions by introducing a time axis.  This has the effect of
`unwinding" the trajectories by stretching them out in the
$t$-direction: closed trajectories become endless coils, equilibrium
points become straight lines parallel to the $t$-axis, and so on.

The analytic tools we introduced to find and explore the nature of
equilibrium points in two-dimensional systems carry over to three
dimensions.  In particular, in investigating the nature of an
equilibrium point analytically, we first localize the system at the
equilibrium point and linearize.  The behavior of the linearized
system can then be explored using analogues of the techniques
introduced in the previous section (or using simple linear algebra).

There are also, of course, limit cycles, which can be attractors,
repellors of a mix of the two (attracting, for example, all
trajectories on a plane, but repelling all trajectories off the plane)
in three-dimensional systems.  As in the case of dynamical systems in
the plane, attracting limit cycles in a three-dimensional system
signal stable periodic behavior.
  
\vspace*{2.5in}
\special{psfile=../ps/calc2/cycle.ps}

\enlargethispage*{20pt}
However, more complicated types of periodic behavior are possible in
the three-dimensional case: we could, for example, have an attracting
torus in the state space.
  
\vspace*{2.5in}
\special{psfile=../ps/calc2/torus.ps}

In this case, the behavior of the states does not settle down to
periodic behavior, but a behavior which is approximately periodic
(often called {\bf quasi-periodic})\index{quasi-periodic}.  In the
plane, there is a well studied phenomenon called Hopf
bifurcation\index{Hopf bifurcation} in which changing the parameters
in a dynamical system can cause an attracting fixed point to become a
repellor surrounded by a stable attracting limit cycle.  Such
dynamical systems arise in modelling situations in which a state
begins to oscillate.  In three dimensions we also sees the same sort
of phenomenon in which an attractor can give birth to an attracting
limit cycle.  However, there are also three-dimensional systems in
which varying the parameters results in an attracting limit cycle
becoming a repelling limit cycle enclosed by an attracting torus (this
is also called Hopf bifurcation and is frequently encountered in
applications).

These sorts of behavior are relatively straightforward generalizations
of behavior in the plane.  At the turn of the century, Poincar\'e
realized that simple three-dimensional systems could have exceedingly
complicated trajectories which exhibit behavior totally unlike any
two-dimensional trajectory. Discoveries in the last three decades have
made it clear that qualitatively new types of {\it attractors} (not
just trajectories) can exist in three-dimensional systems with even
very simple equations.  The most famous such attractor was discovered
by a meteorologist, Edward Lorenz\index{Lorenz, Edward}, in the course
of using dynamical systems to model weather patterns. He discovered a
class of simple systems with an attractor which corresponded to
behavior which was in no sense periodic.  An example of such a system
is
\begin{align*}
x' &= -3x - 3y, \\
y' &= -xz + 30x - y,  \\
z' &= xy - z.
\end{align*}
All trajectories of the system entered a bounded region of the state
space and tended towards a clearly defined geometrical object
(resembling a butterfly). But along the attractor, nearby points
followed trajectories which rapidly diverged from one another. Below,
we have sketched two views of a trajectory beginning at (0,1,0) of the
system above.
 
\vspace*{2.6in}
\special{psfile=../ps/calc2/lor.ps}


\newpage

As Lorenz noted in the paper describing his discovery (Deterministic
Nonperiodic Flow, {\em J. Atmos. Sci.}, {\bf 20} 130 (1963)), this
divergence of trajectories along an attractor has astonishing
practical implications.  It means that that the trajectories of nearby
points in state space could (and would) wind up following very
different paths along the attractor.  Since we never know initial
conditions exactly (and even if we did, a computer truncates decimal
expansions of the coordinates of any point, effectively replacing the
point with a nearby point), this means that long-term predictions
using a model possessing such an attractor are impossible.  In other
words, although the future is completely determined by a dynamical
system given an initial state, it is unknowable in systems of the sort
discovered by Lorenz, because initial states are never known exactly
in practice. Such systems are called {\bf chaotic}\index{chaotic
  system} and attractors which are not points, limit cycles or tori
are called {\bf strange attractors}\index{strange attractor}.  These
systems have been intensively studied in the last thirty years and are
still far from completely understood. Chaotic systems have been used
to attempt to model a wide variety of real situations which exhibit
unpredictable behavior: business cycles, turbulence, heart attacks,
etc.  Although fascinating and philosophically provocative, most of
this work is still very speculative and has yet to prove of practical
value.

Systems involving more than three variables can still be treated
geometrically: we think of the space of states as a higher
dimensional space (one dimension for each quantity) and the dynamical
system as defining a vector field on the state space.  Of course, we
cannot visualize such spaces directly, but the geometrical insight we
gain in dimensions two and three very frequently allows us to handle
such systems.

\subsection{Exercises}

In the next two exercises, we look at some three-dimensional systems
which arise in ecology.  These questions are challenging and you will
probably find it helpful to work them out in a group.

\numa Consider a system consisting of three species: giant carnivorous
reptiles, vegetarian mammals, and plants.  Suppose that the
populations of these are given by $x$, $y$ and $z$ respectively.  The
reptiles eat the mammals, the mammals eat the plants, and the plants
compete among themselves.  Explain why the following system is
consistent with these hypotheses:
\begin{align*}
x' &= -.2 x + .0001 x y, \\
y' &= -.05 y - .001 x y + .000001 y z, \\
z' &= z - .00001 z^2 - .0001 y z.
\end{align*}

\sub Find all equilibrium points of the system.  There are five, one
of which is physically impossible.  Describe the significance of the
other four.

\sub The most interesting equilibrium is the one in which all three
species are present.  Localize the system at this equilibrium, using
local variables $u$, $v$, and $w$. Linearize.  Show that the
linearized system has the form
\begin{alignat*}{2}
u' &=     & \hphantom{-}.003v, &         \\
v' &= -2u &       &+ .002w,  \\
w' &=     & -8v   &- .8w.
\end{alignat*}
Can you determine whether the equilibrium is an attractor?  This is a
hard question---it {\em is\/} an attractor.  One way to show this is a
generalization of the technique we used in the preceding section to
examine the distance of points on a trajectory from the origin over
time.  For the current problem we use a {\bf generalized distance
  function}
\[
D(t) = 8\cdot 10^6u^2+12000v^2+3w^2.
\]
Show, using arguments like those we used when we looked at ordinary
distance, that as we move along a trajectory, the value of $D$ must
decrease.  Hence conclude that the equilibrium point must be an
attractor.

\num The system of equations 
\begin{align*}
x' &= x - .001 x^2 + .002 x y -.1 x z, \\
y' &= y - .01 y ^2 + .001 x y, \\
z' &= -z + .001 x z.
\end{align*}
arises in a general family of models proposed in 1980 by Heithaus,
Culver, and Beattie (``Models of Some Ant-Plant Mutualisms,'' {\em
  American Naturalist}, {\bf 116} (1980) pp. 347-361) for
investigating the interactions three species: violets, ants, and mice.
Violets produce seeds with density $x$ (per square meter, say).  The
ants take some of the seeds and use the seed covering for food.  But
they leave the remainder, which is still a perfectly good seed, in
their refuse piles, which happily turn out to be good sites for
germination.  The ants have density~$y$.  Finally, the seeds are also
taken by mice, who use the whole seed for food (destroying both the
cover and the seed within). The mice have density~$z$.  \sub Explain
why these equations are consistent with the hypotheses we made on the
interactions between the violets, ants and mice.

\sub Find all equilibrium points for the system.  Don't forget the
points where one or more of the variables equals 0.
\begin{align*}
x' &= x - .001 x^2 + .002 x y -.1 x z, \\
y' &= y - .01 y ^2 + .001 x y, \\
z' &= -z + .001 x z.
\end{align*}

\sub Localize the model at each of these equilibria, using local
coordinates $u$, $v$, and $w$ as before, and linearize.

\sub In the case of the equilibrium point (1000, 200, 4) the local
linearization is
\begin{align*}
u' &= -u + 2 v - 100 w, \\
v' &= .2 u - 2 v, \\
w' &= .004 u.
\end{align*}
As in the preceding problem, show that this point is an attractor by
examining the generalized distance function
$R(t)=u(t)^2+v(t)^2+25000w(t)^2$ and showing that the value of $R$
decreases as you move along a trajectory.

%%%%%%%%%%%%%%%%%%%%%%%%%%%%%%%%%%%%%%%%%%%%%%%%%%%%%%%%%%%%%%%%%%%%%%%%%
%                                                                       %
%                                Section 6                              % 
%                                                                       %
%%%%%%%%%%%%%%%%%%%%%%%%%%%%%%%%%%%%%%%%%%%%%%%%%%%%%%%%%%%%%%%%%%%%%%%%%

\section{Chapter Summary}

\subsection{The Main Ideas}

\begin{itemize}

\item A dynamical system can be viewed as a geometrical object. The
  possible values of the dependent variables are then the coordinates
  of a point---called a {\bf state}.  The set of all possible points
  is called the {\bf state space} for the system.

\item The differential equations become a rule assigning a {\bf
  velocity vector} to each state.  Thought of in this way, the
  equations are called a {\bf vector field}.

\item Solutions to the differential equations correspond to {\bf
  trajectories} in the state space.  At every point a trajectory is
  tangent to the corresponding velocity vector, and is changing at the
  rate given by the length of the vector.  The set of all possible
  trajectories is called the {\bf phase portrait} of the system.

\item {\bf Equilibrium points} are points where the velocity vector is
  0. An equilibrium point is a trajectory consisting of a single
  point.  A dynamical system is conveniently analyzed by examining the
  nature of its equilibrium points---whether they are {\bf
    attractors}, {\bf repellors}, {\bf saddle points}, or {\bf
    centers}.

\item To study the nature of an equilibrium point it is helpful to
  look at the {\bf local linearization} of the vector field near the
  point.

\item Determining whether {\bf fixed-line trajectories} exist is a
  crucial part of analyzing the nature of an equilibrium point.

\item In addition to equilibrium points, dynamical systems in two
  dimensions may also have {\bf limit cycles} that shape the
  asymptotic behavior of the system.

\item In higher dimensional state spaces, there are not only the
  obvious extensions of point attractors and limit cycles, but it is
  possible to have an {\bf attracting torus} as well.  There are even
  more complicated attracting objects called {\bf strange attractors}.

\end{itemize}

\subsection{Expectations}

\begin{itemize}

\item You should be able to describe the assumptions embodied in a
  particular dynamical system modeling the interaction between two (or
  three) species and evaluate whether the assumptions seem reasonable.

\item For a dynamical system with two dependent variables, you should
  be able to determine the regions where each variable is zero or has
  a constant sign, find equilibrium points, sketch representative
  vectors of the vector field, and draw trajectories that are
  consistent with this information.

\item You should be able to determine whether a linear system of
  differential equations with two dependent variables has
  \textbf{fixed-line} trajectories---that is, trajectories that are
  straight lines going directly toward or directly away from an
  equilibrium point.

\item You should be able to \textbf{localize} and \textbf{linearize} a
  dynamical system in two variables to explore its behavior near an
  equilibrium point.

\item You should be able to recognize the five generic types of
  equilibrium points: attracting and repelling \textbf{nodes},
  attracting and repelling \textbf{spirals}, and \textbf{saddle
    points}.

\item Using a differential equation solver, you should be able to
  recognize when a dnamical system has a \textbf{limit cycle}.

\item You should be able to analyze a dynamical system with three
  dependent variables.

\end{itemize}

\cleardoublepage

